\label{sec:conclusion}

This paper presents a spatial contagion sandbox that extends the DevAll multi-agent platform, demonstrating a novel integration of cellular automata principles, agent-based modeling, and large language model-powered systems. The artifact contributes a zero-code, configuration-driven approach to epidemiological simulation that addresses a critical accessibility gap in existing platforms. By implementing contagion dynamics through local state transitions, proximity-based transmission, and environmental contamination, the sandbox enables non-programmers---educators, healthcare professionals, and policy analysts---to rapidly prototype and explore disease transmission scenarios without writing code.

The evaluation reveals that simple local transmission rules, inspired by Conway's Game of Life, produce emergent macro-level patterns analogous to compartmental epidemiological models. Observations include spontaneous quarantine behaviors where agents avoid infected individuals through collision resolution, contamination trail dynamics creating persistent transmission reservoirs, and recovery wave patterns propagating through spatial populations. These emergent phenomena validate the cellular automata approach to contagion modeling and demonstrate that complex epidemiological dynamics can arise from composable, local interaction rules without centralized coordination.

The DevAll contagion sandbox is released as open-source software under the MIT License, available at \url{https://github.com/ChatDev-DevAll/multi-agent-simulation}. The platform welcomes community contributions across three dimensions: (1) new transmission scenarios modeling diverse epidemiological contexts (airborne diseases, vector-borne transmission, multi-strain dynamics), (2) visualization extensions including real-time charts, geographic overlays, and statistical dashboards, and (3) pedagogical integrations with educational curricula for K-12 STEM programs and healthcare professional training. By democratizing access to epidemiological simulation tools, this research aims to empower diverse stakeholders in pandemic preparedness, public health education, and intervention design.
